\Needspace{14\baselineskip}
\Article{30}

Another method for determining the sum $W$ in general follows from what we set out in Articles $22.\ 24$.\quad Put $\cos\omega+i\sin\omega=\rho$, and

\[
\begin{aligned}
\rho^{\frac{nn}{aa}}=\alpha,\quad\rho^{\frac{nn}{bb}}=\beta,\quad\rho^{\frac{nn}{cc}}=\gamma\;\text{etc.}
\end{aligned}
\]

so that $r=\rho^k$, $A=\alpha^k$, $B=\beta^k$, $C=\gamma^k$ etc.\quad Then

\[
\begin{aligned}
1+\rho+\rho^4+\rho^9+\text{etc.}+\rho^{(n-1)^2}
\end{aligned}
\]

is the product of the factors

\[
\begin{aligned}
&1+\alpha+\alpha^4+\alpha^9+\text{etc.}+\alpha^{(a-1)^2}\\
&1+\beta+\beta^4+\beta^9+\text{etc.}+\beta^{(b-1)^2}\\
&1+\gamma+\gamma^4+\gamma^9+\text{etc.}+\gamma^{(c-1)^2},\;\text{etc.}
\end{aligned}
\]

\SourceLeaf{50}

\Needspace{12\baselineskip}
and therefore $W$ is the product of the factors

\[
\begin{aligned}
w&=1+\rho+\rho^4+\rho^9+\text{etc.}+\rho^{(n-1)^2}\\[2pt]
\mathfrak A&=\frac{1+A+A^4+A^9+\text{etc.}+A^{(a-1)^2}}{1+\alpha+\alpha^4+\alpha^9+\text{etc.}+\alpha^{(a-1)^2}}\\[2pt]
\mathfrak B&=\frac{1+B+B^4+B^9+\text{etc.}+B^{(b-1)^2}}{1+\beta+\beta^4+\beta^9+\text{etc.}+\beta^{(b-1)^2}}\\[2pt]
\mathfrak C&=\frac{1+C+C^4+C^9+\text{etc.}+C^{(c-1)^2}}{1+\gamma+\gamma^4+\gamma^9+\text{etc.}+\gamma^{(c-1)^2}}
\end{aligned}
\]

The first factor $w$ has already been determined by the investigations given above (Article $19$); the remaining factors $\mathfrak A,\mathfrak B,\mathfrak C$ etc. follow from the formulas of Articles $22.\ 24$, which we collect here again to have everything together\begingroup\renewcommand{\thefootnote}{*)}\footnote{Clearly, what were $k$ and $h$ there are here $\frac{n}{a}$ and $k$ for the second factor, $\frac{n}{b}$ and $k$ for the third factor, etc.}\endgroup.\quad Twelve cases must be distinguished here, namely

I.\quad If $a$ is an (odd) prime number $=p$, or a power of such a number with an odd exponent, and $k$ is a quadratic residue of $p$, the corresponding factor is $\mathfrak A=+1$.

II.\quad If the other conditions remain unchanged and $k$ is a quadratic nonresidue of $p$, then $\mathfrak A=-1$.

III.\quad If $a$ is the square of an odd prime number, or a higher power of it with an even exponent, then $\mathfrak A=+1$.

IV.\quad If $a$ is $=4$, or a higher power of two with an even exponent, and $k$ is of the form $4\mu+1$, then $\mathfrak A=+1$.

V.\quad If, with the other conditions as in IV, $k$ is of the form $4\mu+3$ and $\frac{n}{a}$ of the form $4\mu+1$, then $\mathfrak A=-i$.

VI.\quad If, with the other conditions as in IV, $k$ is of the form $4\mu+3$ and $\frac{n}{a}$ of the form $4\mu+3$, then $\mathfrak A=+i$.

VII.\quad If $a$ is $=8$, or a higher power of two with an odd exponent, and $k$ is of the form $8\mu+1$, then $\mathfrak A=+1$.

VIII.\quad If, with the other conditions as in VII, $k$ is of the form $8\mu+5$, then $\mathfrak A=-1$.

IX.\quad If, with the other conditions as in VII, $k$ is of the form $8\mu+3$ and $\frac{n}{a}$ of the form $4\mu+1$, then $\mathfrak A=+i$.

\SourceLeaf{51}

X.\quad If, with the other conditions as in VII, $k$ is of the form $8\mu+3$ and $\frac{n}{a}$ of the form $4\mu+3$, then $\mathfrak A=-i$.

XI.\quad If, with the other conditions as in VII, $k$ is of the form $8\mu+7$ and $\frac{n}{a}$ of the form $4\mu+1$, then $\mathfrak A=-i$.

XII.\quad If, with the other conditions as in VII, $k$ is of the form $8\mu+7$ and $\frac{n}{a}$ of the form $4\mu+3$, then $\mathfrak A=+i$.

We omit the case where $a=2$; here $\mathfrak A$ would indeed be $=\frac00$, or indeterminate, but then we always have $W=0$.

The remaining factors $\mathfrak B,\mathfrak C$ etc. depend on $b,c$ etc. in the same way as $\mathfrak A$ depends on $a$, insofar as these enter into their determination.

\Needspace{14\baselineskip}
\Article{31}

By this second method, the first example of Article $29$ runs as follows:

{\small
\[
\begin{aligned}
&\text{The factor }w\text{ is }=(1+i)\surd2520\\
&\text{For }a=8\text{ the corresponding factor }\mathfrak A\text{ is, by case VIII, }=-1\\
&\text{For the second factor of }n\text{, }9\text{, the corresponding factor is }+1\;\text{(by case III.)}\\
&\text{For the factor }5\text{ the corresponding factor is }-1\;\text{(by case II.)}\\
&\text{For the factor }7\text{ the corresponding factor is }-1\;\text{(by case II.)}
\end{aligned}
\]
}

This gives the product $W=(-1-i)\surd2520$, as in Article $29$.

\Needspace{18\baselineskip}
\Article{32}

Since the value of $W$ can be determined by \textit{two} methods, one of which rests on the relations of the numbers $\frac{nk}{a}$, $\frac{nk}{b}$, $\frac{nk}{c}$ etc. to the numbers $a,b,c$ etc., while the other depends on the relations of $k$ to the numbers $a,b,c$ etc., there must be a conditional connection among all these relations, so that any one can be determined from the others.\quad Suppose that all the numbers $a,b,c$ etc. are odd prime numbers, and take $k$ $=1$. Divide the factors $a,b,c$ etc. into two classes: one contains those of the form $4\mu+1$, denoted by $p,p',p''$ etc.; the other consists of those of the form $4\mu+3$, denoted by $q,q',q''$ etc. We shall denote the number of the latter by $m$. With this done, we first observe that $n$ is of the form $4\mu+1$ if $m$ is even (this also includes the case where there are no factors at all in the second class, or $m=0$), whereas $n$ is of the form $4\mu+3$ if $m$ is odd.\quad Now the determination\SourceLeaf{52} of $W$ by the first method is carried out as follows.\quad Let the numbers $P,P'\;P''$ etc., $Q,Q',Q''$ etc. depend on the relations of the numbers $\frac{n}{p}$, $\frac{n}{p'}$, $\frac{n}{p''}$ etc., $\frac{n}{q}$, $\frac{n}{q'}$, $\frac{n}{q''}$ etc. to the numbers $p,p',p''$ etc., $q,q',q''$ etc., respectively, in such a way that we put

\[
\begin{aligned}
P&=+1,\;\text{if }\frac{n}{p}\text{ is a quadratic residue of }p\\
P&=-1,\;\text{if }\frac{n}{p}\text{ is a quadratic nonresidue of }p
\end{aligned}
\]

and similarly for the others.\quad Then $W$ will be the product of the factors $P\surd p$, $P'\surd p'$, $P''\surd p''$ etc. $iQ\surd q$, $iQ'\surd q'$, $iQ''\surd q''$ etc., and therefore

\[
\begin{aligned}
W=PP'P''\ldots QQ'Q''\ldots i^m\surd n
\end{aligned}
\]

\Needspace{10\baselineskip}
By the second method, or rather directly from the rules of Article $19$, we have

\[
\begin{aligned}
W&=+\surd n,\;\text{if }n\text{ is of the form }4\mu+1,\\&\qquad\text{or equivalently, if }m\text{ is even}\\
W&=+i\surd n,\;\text{if }n\text{ is of the form }4\mu+3,\\&\qquad\text{or if }m\text{ is odd.}
\end{aligned}
\]

Both cases can be combined in the following formula:

\[
\begin{aligned}
W=i^{mm}\surd n
\end{aligned}
\]

It therefore follows that

\[
\begin{aligned}
PP'P''\ldots QQ'Q''\ldots=i^{mm-m}
\end{aligned}
\]

But $i^{mm-m}$ is $=1$ whenever $m$ is of the form $4\mu$ or $4\mu+1$, and $=-1$ whenever $m$ is of the form $4\mu+2$ or $4\mu+3$. From this we deduce the following most elegant

\textsc{Theorem.}\quad \textit{Let $a,b,c$ etc. denote distinct positive odd prime numbers, and put their product $=n$. Let $m$ of them be of the form $4\mu+3$, and the rest of the form $4\mu+1$. The number of those among $a,b,c$ etc. for which $\frac{n}{a}$, $\frac{n}{b}$, $\frac{n}{c}$ etc., respectively, are nonresidues will be even whenever $m$ is of the form $4\mu$ or $4\mu+1$, but odd whenever $m$ is of the form $4\mu+2$ or $4\mu+3$.}

Thus, for example, putting $a=3$, $b=5$, $c=7$, $d=11$, we have three numbers of the form $4\mu+3$, namely $3$, $7$ and $11$; and $5\pdot7\pdot11\,R3$; $3\pdot7\pdot11\,R5$; $3\pdot5\pdot11\,R7$; $3\pdot5\pdot7\,N11$, so that only $\frac{n}{d}$ is a nonresidue of $d$.

\Needspace{12\baselineskip}
\Article{33}

The celebrated \textit{fundamental theorem} about quadratic residues is nothing other than a special case of the theorem just developed.\quad Specifically, by limiting the number\SourceLeaf{53} of the numbers $a,b,c$ etc. to \textit{two}, it is clear that, if just one of them, or neither, is of the form $4\mu+3$, we must have either both $aRb$, $bRa$, or both $aNb$, $bNa$; but if both are of the form $4\mu+3$, one must be a nonresidue of the other, and the latter a residue of the former.\quad Here, then, is a \textit{fourth} proof of this most important theorem, whose first and second proofs we gave in the Disquisitiones Arithmeticae, and the third recently in a separate paper (\textit{Commentt. T. XVI}); we shall later set out two others, again resting on entirely different principles.\quad It is certainly most remarkable that this most beautiful theorem, which at first so stubbornly eluded every attempt, could later be reached by so many paths, worlds apart from one another.

\Needspace{14\baselineskip}
\Article{34}

The remaining theorems, which form a kind of supplement to the fundamental theorem, namely those by which we distinguish the prime numbers for which $-1,+2$ and $-2$ are residues or nonresidues, can also be derived from the same principles. We shall begin with the residue $+2$.

Putting $n=8a$, with $a$ a prime number, and $k=1$, the method of Article $28$ makes $W$ the product of two factors. One of them will be $+\surd a$, or $+i\surd a$, if $8$, or equivalently $2$, is a quadratic residue of $a$; but $-\surd a$ or $-i\surd a$ if $2$ is a nonresidue of $a$.\quad The second factor is

\[
\begin{aligned}
&(1+i)\surd8,\quad\text{if }a\text{ is of the form }8\mu+1\\
&(-1+i)\surd8,\quad\text{if }a\text{ is of the form }8\mu+3\\
&(-1-i)\surd8,\quad\text{if }a\text{ is of the form }8\mu+5\\
&(1-i)\surd8,\quad\text{if }a\text{ is of the form }8\mu+7
\end{aligned}
\]

But by Article $18$ we always have $W=(1+i)\surd n$; dividing this value by the four values of the second factor shows that the first factor must be

\[
\begin{aligned}
&+\surd a,\quad\text{if }a\text{ is of the form }8\mu+1\\
&-i\surd a,\quad\text{if }a\text{ is of the form }8\mu+3\\
&-\surd a,\quad\text{if }a\text{ is of the form }8\mu+5\\
&+i\surd a,\quad\text{if }a\text{ is of the form }8\mu+7
\end{aligned}
\]

It follows immediately that $2$ must be a residue of $a$ in the first and fourth cases, but a nonresidue in the second and third.

\SourceLeaf{54}

\Needspace{14\baselineskip}
\Article{35}

The prime numbers for which $-1$ is a residue or a nonresidue are easily distinguished with the aid of the following theorem, which is also quite remarkable in itself.

\textsc{Theorem.}\quad \textit{The product of the two factors}

\[
\begin{aligned}
W'&=1+r^{-1}+r^{-4}+\text{etc.}+r^{-(n-1)^2}\\
W&=1+r+r^4+\text{etc.}+r^{(n-1)^2}
\end{aligned}
\]

\textit{is $=n$ if $n$ is odd; or $=0$ if $n$ is twice an odd number; or $=2n$ if $n$ is divisible by four.}

\Needspace{12\baselineskip}
\textit{Proof.}\quad Since we evidently have

\[
\begin{aligned}
W&=r+r^4+r^9+\text{etc.}+r^{nn}\\
&=r^4+r^9+\text{etc.}+r^{(n+1)^2}\\
&=r^9+\text{etc.}+r^{(n+2)^2}\quad\text{etc.}
\end{aligned}
\]

the product $WW'$ can also be written as follows:

\[
\begin{aligned}
&1+r+r^4+r^9+\text{etc.}+r^{(n-1)^2}\\
+{}&r^{-1}(r+r^4+r^9+r^{16}+\text{etc.}+r^{nn})\\
+{}&r^{-4}(r^4+r^9+r^{16}+r^{25}+\text{etc.}+r^{(n+1)^2})\\
+{}&r^{-9}(r^9+r^{16}+r^{25}+r^{36}+\text{etc.}+r^{(n+2)^2})\\
&\text{etc.}\\
+{}&r^{-(n-1)^2}(r^{(n-1)^2}+r^{nn}+r^{(n+1)^2}\\&\qquad+r^{(n+2)^2}+\text{etc.}+r^{(2n-2)^2})
\end{aligned}
\]

Summing this expression vertically gives

\[
\begin{aligned}
&n\\
+{}&r(1+rr+r^4+r^6+\text{etc.}+r^{2n-2})\\
+{}&r^4(1+r^4+r^8+r^{12}+\text{etc.}+r^{4n-4})\\
+{}&r^9(1+r^6+r^{12}+r^{18}+\text{etc.}+r^{6n-6})\\
+{}&\text{etc.}\\
+{}&r^{(n-1)^2}(1+r^{2n-2}+r^{4n-4}\\&\qquad+r^{6n-6}+\text{etc.}+r^{2(n-1)^2})
\end{aligned}
\]

Now, if $n$ is odd, each part of this expression, except the first, $n$, will be $=0$; for the second is evidently $\frac{r(1-r^{2n})}{1-rr}$, the third $\frac{r^4(1-r^{4n})}{1-r^4}$ etc. Whenever $n$ is even, however, we must also except the part

\SourceLeaf{55}

\[
\begin{aligned}
r^{\frac14nn}(1+r^n+r^{2n}+r^{3n}+\text{etc.}+r^{nn-n})
\end{aligned}
\]

which is $=nr^{\frac14nn}$. In the first case, therefore, $WW'=n$, but in the second $=n+nr^{\frac14nn}$; now $r^{\frac14nn}$ is $=+1$ if $n$ is divisible by four, and then we obtain $WW'=2n$; on the other hand, $r^{\frac14nn}=-1$ if $n$ is twice an odd number, in which case $WW'=0$. Q. E. D.

\Needspace{10\baselineskip}
\Article{36}

Article 22 establishes that, if $n$ is an odd prime number, $\frac{W'}{W}$ is $=+1$ or $=-1$ according as $-1$ is a residue or a nonresidue of $n$. Hence in the first case we must have $W^2=+n$, and in the second $W^2=-n$; by Article 13 we therefore conclude that the first case can occur only when $n$ is of the form $4\mu+1$, and the second when $n$ is of the form $4\mu+3$.

Finally, combining the conditions found for the residues $+2$ and $-1$ immediately shows that $-2$ is a residue of every prime number of the form $8\mu+1$ or $8\mu+3$, and a nonresidue of every prime number of the form $8\mu+5$ or $8\mu+7$.

\BookRule{42mm}
